% Shared technical body for the two individual D1 Part 2 PDFs.
\usepackage[utf8]{inputenc}
\usepackage[T1]{fontenc}
\usepackage{mathpazo}
\usepackage[scaled=0.92]{helvet}
\usepackage[varqu,varl]{zi4}
\usepackage[english]{babel}
\usepackage{microtype}
\usepackage[a4paper,margin=2.15cm,top=1.8cm,bottom=2.0cm]{geometry}
\usepackage[table]{xcolor}
\definecolor{ansblue}{HTML}{1E3A8A}
\definecolor{anscodebg}{HTML}{F9FAFB}
\definecolor{anscaption}{HTML}{6B7280}
\definecolor{anskeyword}{HTML}{2563EB}
\usepackage[colorlinks=true,linkcolor=ansblue,urlcolor=ansblue,citecolor=ansblue,pdfborder={0 0 0}]{hyperref}
\usepackage[explicit]{titlesec}
\titleformat{\section}{\large\bfseries\sffamily}{\thesection}{0.65em}{#1}
\titleformat{\subsection}{\normalsize\bfseries\sffamily}{\thesubsection}{0.65em}{#1}
\titlespacing*{\section}{0pt}{10pt}{4pt}
\titlespacing*{\subsection}{0pt}{6pt}{2pt}
\setlength{\parindent}{0pt}
\setlength{\parskip}{4pt}
\linespread{1.04}
\setcounter{secnumdepth}{0}
\usepackage{array}
\usepackage{tabularx}
\usepackage{booktabs}
\usepackage{ragged2e}
\newcolumntype{Y}{>{\RaggedRight\arraybackslash}X}
\usepackage{enumitem}
\setlist[itemize]{leftmargin=1.1em,itemsep=1pt,topsep=2pt}
\usepackage{listings}
\lstdefinelanguage{ANS}{
  morekeywords={let,function,return,for,in,if,elif,else,print,element,compound,reaction,reactants,products,formula,name,atomic_mass,valence,symbol,check,balanced,energy,exothermic,endothermic,mol},
  sensitive=true,
  morecomment=[l]{//},
  morestring=[b]",
  morestring=[b]'
}
\lstdefinestyle{ans}{
  language=ANS,
  basicstyle=\ttfamily\footnotesize,
  keywordstyle=\bfseries\color{anskeyword},
  commentstyle=\itshape\color{anscaption},
  stringstyle=\color{anskeyword!80!black},
  numbers=none,
  frame=none,
  backgroundcolor=\color{anscodebg},
  xleftmargin=5pt,
  xrightmargin=5pt,
  aboveskip=5pt,
  belowskip=5pt,
  showstringspaces=false,
  columns=fullflexible,
  keepspaces=true,
  breaklines=true,
  upquote=true
}
\lstdefinestyle{ebnf}{
  basicstyle=\ttfamily\scriptsize,
  numbers=none,
  frame=none,
  backgroundcolor=\color{white},
  columns=fullflexible,
  keepspaces=true,
  breaklines=true,
  literate={->}{{$\rightarrow$}}2
}
\newcommand{\code}[1]{\texttt{#1}}
\usepackage{fancyhdr}
\pagestyle{fancy}
\fancyhf{}
\renewcommand{\headrulewidth}{0pt}
\fancyfoot[L]{\footnotesize\color{anscaption}CSE 341: ANS Design Specification (D1, Part 2) \(\cdot\) \DocAuthor}
\fancyfoot[R]{\footnotesize\color{anscaption}Page \thepage}
\hypersetup{
  pdftitle={ANS Design Specification (D1, Part 2)},
  pdfauthor={\DocAuthor},
  pdfsubject={CSE 341 Project Part 2},
  pdfkeywords={DSL, chemistry, type system, interpreter, Sebesta}
}

\begin{document}
\begin{center}
  {\footnotesize\color{anscaption}CSE 341: Concepts of Programming Languages\par}
  \vspace{3pt}
  {\huge\bfseries\sffamily ANS: Alchemist's Notation Script\par}
  \vspace{4pt}
  {\large Design Specification (D1), Part 2\par}
  \vspace{5pt}
  {\footnotesize\color{anscaption}Student: \StudentID\ \StudentName\ \(\cdot\) Spring 2026 \(\cdot\) Gebze Technical University\par}
\end{center}

\section{4.1 Language Overview}
\textbf{Name and domain.} ANS, Alchemist's Notation Script, is a small chemistry DSL for declaring elements, compounds, and reactions, then running simple stoichiometric calculations over them. Its surface syntax keeps chemistry visible: users write \code{compound water \{ formula = H2O \}}, \code{CH4 + 2*O2}, and \code{check = balanced} instead of hiding those concepts in strings or generic arrays.

\begin{lstlisting}[style=ans]
element H { atomic_mass = 1.008,  valence = 1, symbol = "H" };
element O { atomic_mass = 15.999, valence = 2, symbol = "O" };
compound water { formula = H2O, name = "Water" };

function molar_mass(compound c) {
    let formula = c.formula;
    let mass = 0.0;
    for (let element in formula) {
        mass += element.atomic_mass * formula[element];
    }
    return mass;
};

if (molar_mass(water) > 18.0) {
    print "Water molar mass is above 18.";
}
\end{lstlisting}

Using Sebesta's Chapter 1 criteria, ANS prioritizes readability, writability, and reliability. Readability comes from domain vocabulary and formula literals. Writability comes from making compounds, reactions, units, and formula traversal first-class. Reliability is the main priority: formula symbols, typed parameters, and balanced reactions are checked before execution. The accepted trade-off is cost. ANS is not a general-purpose language, and the implementation must know about atom counts, reaction sides, and chemistry-specific lookup errors.

\newpage
\section{4.2 Lexical Structure}
ANS has identifiers, literals, hard keywords, contextual chemistry words, operators, separators, comments, and whitespace. Identifiers match \code{[a-z\_][A-Za-z0-9\_]*}; formula literals match uppercase chemical notation such as \code{H}, \code{O2}, and \code{CH4}. This case split keeps \code{water} and \code{H2O} separate without needing a stateful lexer.

{\footnotesize
\begin{tabularx}{\linewidth}{@{}p{2.8cm}p{4.0cm}Y@{}}
\toprule
Category & Examples & Rule \\
\midrule
Identifiers & \code{water}, \code{mass} & lowercase or underscore start \\
Numbers & \code{20}, \code{1.008}, \code{6.02e23} & integer, decimal, or scientific notation \\
Strings/chars & \code{"Water"}, \code{'H'} & quoted literals with simple escapes \\
Formula literals & \code{H}, \code{H2O}, \code{CH4} & repeated element-symbol groups \\
Unit literals & \code{g/mol}, \code{kg/mol}, \code{kJ/mol} & longest match before slash division \\
Single-word units & \code{mol}, \code{g}, \code{kg}, \code{kJ} & identifiers promoted by parser after a number \\
Keywords & \code{let}, \code{function}, \code{return}, \code{for}, \code{if}, \code{print} & hard reserved words \\
Operators/separators & \code{+ - * / == != <= >= += = \{ \} ( ) [ ] , ; .} & maximal munch \\
Comments & \code{// comment} & skipped to newline \\
\bottomrule
\end{tabularx}}

\newpage
\section{4.3 Syntax}
\begin{lstlisting}[style=ebnf]
<program> -> { <top_level_statement> } EOF
<top_level_statement> -> <element_declaration> ";"
                       | <compound_declaration> ";"
                       | <reaction_declaration> ";"
                       | <function_declaration> ";"
                       | <variable_declaration> ";"
                       | <assignment_statement> ";"
                       | <augmented_assignment_statement> ";"
                       | <print_statement> ";"
                       | <expression_statement> ";"
                       | <for_statement> | <if_statement>

<block> -> "{" { <statement> } "}"
<statement> -> <variable_declaration> ";" | <assignment_statement> ";"
             | <augmented_assignment_statement> ";" | <return_statement> ";"
             | <print_statement> ";" | <expression_statement> ";"
             | <for_statement> | <if_statement>
<variable_declaration> -> "let" <identifier> "=" <expression>
<assignment_statement> -> <assignable> "=" <expression>
<augmented_assignment_statement> -> <assignable> ("+=" | "-=" | "*=" | "/=") <expression>
<return_statement> -> "return" [ <expression> ]
<print_statement> -> "print" <expression>

<function_declaration> -> "function" <identifier> "(" [ <parameter_list> ] ")" <block>
<parameter_list> -> <parameter> { "," <parameter> }
<parameter> -> [ <type_name> ] <identifier>
<type_name> -> "element" | "compound" | "reaction"
<for_statement> -> "for" "(" "let" <identifier> "in" <expression> ")" <block>
<if_statement> -> "if" "(" <expression> ")" <block> { <elif_clause> } [ <else_clause> ]
<elif_clause> -> "elif" "(" <expression> ")" <block>
<else_clause> -> "else" <block>

<element_declaration> -> "element" <formula_literal> "{" <element_field_list> "}"
<element_field> -> "atomic_mass" "=" <number_literal> | "valence" "=" <number_literal>
                 | "symbol" "=" <string_literal>
<compound_declaration> -> "compound" <identifier> "{" <compound_field_list> "}"
<compound_field> -> "formula" "=" <formula_literal> | "name" "=" <string_literal>

<reaction_declaration> -> "reaction" <identifier> "{" <reaction_field_list> "}"
<reaction_field> -> "reactants" "=" <reaction_expression>
                  | "products" "=" <reaction_expression>
                  | "check" "=" "balanced"
                  | "energy" "=" ("exothermic" | "endothermic") <quantity_literal>
<reaction_expression> -> <reaction_term> { "+" <reaction_term> }
<reaction_term> -> <species_ref> | <number_literal> "*" <species_ref>
<species_ref> -> <formula_literal> | <identifier>

<expression> -> <comparison_expression>
<comparison_expression> -> <additive_expression>
                           [ ("==" | "!=" | ">=" | "<=" | ">" | "<") <additive_expression> ]
<additive_expression> -> <multiplicative_expression> { ("+" | "-") <multiplicative_expression> }
<multiplicative_expression> -> <unary_expression> { ("*" | "/") <unary_expression> }
<unary_expression> -> [ "-" ] <postfix_expression>
<postfix_expression> -> <primary_expression> { "." <identifier> | "[" <expression> "]" | "(" [ <argument_list> ] ")" }
<primary_expression> -> <number_or_quantity_literal> | <string_literal> | <char_literal>
                      | <list_literal> | <formula_literal> | <identifier> | "(" <expression> ")"
<argument_list> -> <expression> { "," <expression> }
<number_or_quantity_literal> -> <number_literal> [ <unit> ]
<quantity_literal> -> <number_literal> <unit>
<unit> -> UNIT_LITERAL | "mol" | "g" | "kg" | "kJ"
<list_literal> -> "[" [ <argument_list> ] "]"
<assignable> -> <identifier> { "." <identifier> | "[" <expression> "]" }
\end{lstlisting}

Ambiguity is resolved in the grammar. Arithmetic precedence is encoded by expression levels, all binary operators are left-associative, and reaction expressions use a separate grammar so \code{2*O2} means coefficient times species rather than ordinary multiplication. Braced blocks remove the dangling-\code{else} problem.

\newpage
\section{4.4 Semantics}
ANS uses big-step operational semantics. Let \(\rho\) be the environment, \(S\) the runtime store of declared chemistry objects, and \(\Downarrow\) the evaluation relation.

\textbf{For over a formula.} If expression \(E\) evaluates to a formula with entries \((e_1,n_1)\ldots(e_k,n_k)\), then the loop variable is bound successively to each element object, not to the numeric count. Counts are obtained by indexing the formula.
\[
\frac{\rho,S \vdash E \Downarrow [(e_1,n_1),\ldots,(e_k,n_k)] \quad
      \rho[x \mapsto e_1],S \vdash B \Downarrow S_1 \ldots
      \rho[x \mapsto e_k],S_{k-1} \vdash B \Downarrow S_k}
     {\rho,S \vdash \texttt{for (let x in E) B} \Downarrow S_k}
\]
The expression is evaluated once before iteration. This makes loops stable even if the body mutates unrelated variables.

\textbf{Checked reaction declaration.} A reaction declaration first evaluates each side into a map from compounds to coefficients. Each compound expands to its formula atom counts. If \code{check = balanced}, the declaration succeeds only when every declared element has the same total count on both sides.
\[
\frac{\forall a.\ \sum_{c \in R} coeff_R(c)\cdot atoms(c,a)
      = \sum_{c \in P} coeff_P(c)\cdot atoms(c,a)}
     {\rho,S \vdash \texttt{reaction r \{ reactants=R, products=P, check=balanced \}} \Downarrow S[r \mapsto reaction(R,P)]}
\]
If any element total differs, the semantic checker rejects the program before execution.

\textbf{Function calls.} ANS uses call by value. Arguments are evaluated left to right, copied into a new stack-dynamic activation record, and names inside the function resolve by static scope: local frame, enclosing lexical frame, then global declarations.

\newpage
\section{4.5 Type System}
Primitive types are \code{number}, \code{string}, \code{char}, and \code{bool}. \code{number} covers integer-looking and floating-point numeric literals; comparisons produce \code{bool}. Domain and structured types are \code{element}, \code{compound}, \code{formula}, \code{reaction}, \code{reaction\_side}, \code{quantity}, \code{list<T>}, and function values.

ANS is strongly typed in Sebesta's sense: operations are defined only for compatible operand types and errors are caught before execution when they can be determined statically. Examples: arithmetic requires numbers, \code{if} conditions require bools, \code{formula[element]} returns a number, and \code{reaction\_side[compound]} returns a coefficient. The final examples use homogeneous lists; a mixed list literal is rejected by the Part 2 checker.

ANS does not allow general implicit coercion. There is no automatic string-to-number, compound-to-formula, or quantity-to-number conversion. The only explicit conversion-like operation is the built-in \code{mol(x)}, which extracts the numeric amount from a quantity or returns a numeric coefficient unchanged. This keeps unit handling visible.

For type equivalence, ANS uses name equivalence for domain record types: \code{compound}, \code{formula}, and \code{reaction} are not interchangeable even if their runtime representation contains maps. For \code{list<T>}, equivalence is structural on the element type, so two \code{list<compound>} values are compatible.

Lists are dynamic-length, zero-indexed, and checked at runtime for index bounds. Domain records are implicit heap-dynamic runtime objects managed by the interpreter; ANS exposes no pointers, \code{new}, or \code{delete}.

\newpage
\section{4.6 Names, Binding, Scope, and Lifetime}
Legal identifiers match \code{[a-z\_][A-Za-z0-9\_]*}. ANS is case-sensitive. Uppercase tokens are formula literals, so \code{water} and \code{H2O} cannot collide.

Compile-time bindings include lexemes to tokens, declarations to names, function parameter names to declared domain types, formula symbols to declared elements, compound names to formulas, and checked reactions to atom-count constraints. Runtime bindings include argument values, local variables, loop variables, and computed list or quantity values.

ANS uses static scoping. Name lookup follows the source structure, not the call chain. This matters for reliability: \code{molar\_mass} should not change meaning because a caller has a local variable named \code{water}. Top-level element, compound, reaction, and function declarations have static lifetime. Function parameters, local \code{let} variables, and loop variables are stack-dynamic. Lists and chemistry objects are implicit heap-dynamic but not manually managed by the programmer.

\newpage
\section{4.7 Expressions and Assignment}
{\footnotesize
\begin{tabularx}{\linewidth}{@{}p{3.0cm}p{4.6cm}Y@{}}
\toprule
Level & Operators/forms & Associativity \\
\midrule
Postfix & call, member access, index & left to right \\
Unary & \code{-} & right to left for the single unary operator \\
Multiplicative & \code{* /} & left \\
Additive & \code{+ -} & left \\
Comparison/equality & \code{< <= > >= == !=} & non-chaining in grammar \\
\bottomrule
\end{tabularx}}

ANS currently has comparison operators but no \code{\&\&} or \code{||}, so there are no Boolean short-circuit operators to define. Operand and argument evaluation order is left to right and defined. Assignment is a statement, not an expression, so ANS has no assignment inside conditions or larger expressions.

\newpage
\section{4.8 Design Rationale}
\textbf{Strong typing.} I chose strong typing because ANS is meant to prevent chemistry mistakes, not just calculate numbers. The trade-off is less convenience: users must pass a \code{compound} where a compound is expected and cannot rely on implicit conversion from strings or raw formulas.

\textbf{No broad implicit coercion.} Units are easy to lose in ordinary scripts. By requiring \code{mol(x)} when extracting a numeric amount, the program marks the point where quantity information is intentionally reduced. The trade-off is a slightly more verbose expression.

\textbf{Name equivalence for domain records.} A reaction side, formula, and compound can all be represented internally as maps, but they mean different things. Name equivalence catches accidental interchange. The trade-off is that users sometimes need explicit member access such as \code{c.formula}.

\textbf{Static scoping.} Static scoping makes code readable and supports compile-time checks. Dynamic scoping would make chemistry names depend on the call history, which is not acceptable for a reliability-first DSL. The trade-off is less flexibility for ad-hoc rebinding.

\textbf{Assignment as a statement.} Keeping assignment out of expressions avoids C-style surprises in conditions. The trade-off is less compact code, but ANS values clear calculation steps over terse expression tricks.

\end{document}
